\ifmostrarGabarito
\begin{minted}{python}
import numpy as np

def calcular_matriz_distancias(sensores: np.ndarray, alvos: np.ndarray) -> np.ndarray:
    """
    Calcula a matriz de distâncias euclidianas entre S sensores e A alvos usando broadcasting.
    Retorna uma matriz de dimensões (S, A).
    """
    # Expande dimensões para permitir broadcasting: (S, 1, 2) e (1, A, 2)
    # A diferença produzirá um tensor de coordenadas (S, A, 2)
    diferencas = sensores[:, np.newaxis, :] - alvos[np.newaxis, :, :]
    
    # Eleva ao quadrado, soma no eixo das coordenadas cartesianas (axis=-1) e extrai a raiz
    distancias = np.sqrt(np.sum(diferencas ** 2, axis=-1))
    
    return np.round(distancias, 2)
\end{minted}

\textbf{Explicação:}
O uso de \texttt{np.newaxis} ajusta as formas dos arrays para $(S, 1, 2)$ e $(1, A, 2)$. As regras de \emph{broadcasting} do NumPy replicam virtualmente as dimensões de tamanho 1 sem alocação extra de memória, calculando todas as $S \times A$ combinações de pares simultaneamente em instruções vetoriais ultra-rápidas.
\fi
